\pdfvariable suppressoptionalinfo 611
\documentclass[9pt]{extarticle}
\usepackage[a4paper,left=0.65in,right=0.65in,top=0.55in,bottom=0.55in]{geometry}
\usepackage{fontspec}
\setmainfont{Latin Modern Roman}
\setsansfont{Latin Modern Sans}
\setmonofont{Latin Modern Mono}
\usepackage[english]{babel}
\babelprovide[import]{chinese-simplified}
\babelfont[chinese-simplified]{rm}[Renderer=Harfbuzz,Language=Default]{Droid Sans Fallback}
\usepackage{amsmath,amssymb,booktabs,xurl,titlesec}
\titlespacing*{\section}{0pt}{1.5ex plus .2ex minus .1ex}{.7ex}
\newtheorem{theorem}{Theorem}
\title{Bulgarian Solitaire: Recurrent Necklaces, Every Fixed Iterate,\\
and the Full Finite Koopman Spectrum}
\author{Anonymous structural certificate HCS-C190}
\date{}
\begin{document}
\maketitle
\begin{abstract}
For Bulgarian solitaire on every integer-partition set $\mathcal P(N)$, we
turn Brandt's attributed recurrent classification into a complete periodic
ledger.  The unique decomposition $N=\binom{k}{2}+r$, $0\le r<k$, identifies
recurrent partitions with length-$k$, weight-$r$ binary words, and one
solitaire move becomes rotation.  We derive a closed formula for every
positive-iterate fixed count, then recover all least periods, primitive
cycles, and the finite Artin--Mazur zeta by M\"obius inversion.  The finite
cycle factors also give the full Koopman algebraic spectrum, including the
transient zero multiplicity.  The finite census is a regression oracle, not
the proof of the all-$N$ theorem.
\end{abstract}
\noindent\textbf{Keywords:} Bulgarian solitaire; integer partitions; binary
necklaces; primitive cycles; dynamical zeta; Koopman spectrum.

\begin{otherlanguage}{chinese-simplified}
\begin{center}{\large 中文摘要}\end{center}
本文研究任意整数分拆集合上的保加利亚单人纸牌动力学。利用归属于
\foreignlanguage{english}{Brandt}的循环分拆分类，将循环状态与长度为$k$、重量为$r$
的二进制词对应，其中$N=\binom{k}{2}+r$且$0\le r<k$；一次动力学迭代正好对应循环旋转。
由此得到所有正迭代固定点数的闭式公式，再用\foreignlanguage{english}{M\"obius}反演恢复
全部最小周期、原始循环、有限动力学$\zeta$函数与包含瞬态零重数的完整有限
\foreignlanguage{english}{Koopman}代数谱。有限枚举仅用于回归检验，不替代全参数定理。

\noindent 关键词：保加利亚单人纸牌；整数分拆；二进制项链；原始循环；动力学$\zeta$函数；
\foreignlanguage{english}{Koopman}谱。
\end{otherlanguage}

\section{The noninvertible map and Brandt coordinates}
Write a partition of $N$ as $\lambda=(\lambda_1\ge\cdots\ge\lambda_m>0)$.
The Bulgarian move is
\[
 T_N(\lambda)=\operatorname{sort}\bigl(m,\lambda_1-1,\ldots,
 \lambda_m-1\bigr),
\]
after zero parts are deleted.  This map is generally noninjective.  Every
$N\ge1$ has a unique representation
\[
 N=\binom{k}{2}+r,\qquad k\ge2,\quad 0\le r<k.
\]
Brandt's cyclic-partition theorem~\cite{Brandt82}, treated also by Akin and
Davis~\cite{AkinDavis85}, says that the recurrent set consists exactly of
\[
 \phi(w)=\text{positive parts of }(k-1,k-2,\ldots,0)+w,
 \qquad w\in\{0,1\}^k,\quad |w|=r.
\]
For the right rotation $(\rho w)_i=w_{i-1\bmod k}$, direct subtraction of the
first column gives $T_N\phi=\phi\rho$.  States correspond to words; dynamical
cycles correspond to binary necklaces.

\section{Every fixed iterate and every primitive cycle}
\begin{theorem}\label{thm:ledger}
Let $t\ge1$ and $g=\gcd(k,t)$.  Then
\[
 F_t:=\#\operatorname{Fix}(T_N^t)=
 \begin{cases}
 \displaystyle\binom{g}{rg/k},&(k/g)\mid r,\\[2pt]
 0,&\text{otherwise}.
 \end{cases}
\]
For each $d\mid k$, the number $P_d$ of points of least period $d$ and the
number $C_d$ of $d$-cycles are
\[
 P_d=\sum_{e\mid d}\mu(d/e)F_e,\qquad C_d=P_d/d.
\]
Consequently the zeta function of the full noninvertible finite map is
\begin{equation}\label{eq:zeta}
 \zeta_{T_N}(z)=\exp\!\left(\sum_{t\ge1}\frac{F_t}{t}z^t\right)
 =\prod_{d\mid k}(1-z^d)^{-C_d}.
\end{equation}
\end{theorem}

\paragraph{Proof.}
The index rotation $\rho^t$ has $g$ cycles, each of length $k/g$.  A fixed
binary word is constant on each index cycle, so its weight is a multiple of
$k/g$; choosing the $rg/k$ one-cycles gives the displayed binomial.  A point
fixed by $T_N^t$ is periodic and hence recurrent, so no transient partition
adds to this count.  Divisor M\"obius inversion gives $P_d$, and division by
$d$ gives $C_d$.  Substituting
$F_t=\sum_{d\mid t}dC_d$ into the exponential definition proves
\eqref{eq:zeta}. \hfill$\square$

\section{The full finite Koopman algebraic spectrum}
Let $p(N)=|\mathcal P(N)|$ and let $(U_Nf)(\lambda)=f(T_N\lambda)$ on all
complex functions on $\mathcal P(N)$.  Ordering each functional-graph
component from its transient trees toward its cycle gives
\begin{equation}\label{eq:char}
 \det(\xi I-U_N)=\xi^{p(N)-\binom{k}{r}}
 \prod_{d\mid k}(\xi^d-1)^{C_d},\qquad
 \det(I-zU_N)=\prod_{d\mid k}(1-z^d)^{C_d}=\zeta_{T_N}(z)^{-1}.
\end{equation}
Thus zero has algebraic multiplicity $p(N)-\binom{k}{r}$.  If
$\omega_k=e^{2\pi i/k}$, the nonzero multiplicities are
\[
 \operatorname{mult}(\omega_k^j)
 =\sum_{\substack{d\mid k\\k\mid jd}}C_d,\qquad
 \operatorname{Tr}(U_N^t)=F_t.
\]
Indeed, transient vertices contribute only a zero diagonal block to the
characteristic polynomial, whereas a $d$-cycle contributes $\xi^d-1$ and
every $d$th root once.  Formula~\eqref{eq:char} determines eigenvalues with
algebraic multiplicities, not the nilpotent Jordan block sizes in the
transient sector.

For $N=8=\binom42+2$, the six recurrent words form one $2$-cycle and one
$4$-cycle.  Hence $(F_1,F_2,F_3,F_4)=(0,2,0,6)$,
$\zeta_{T_8}(z)=((1-z^2)(1-z^4))^{-1}$, zero has multiplicity $22-6=16$,
and the four fourth-root multiplicities are $(2,1,2,1)$.  The sentinel is
obtained both from all $22$ partitions of $8$ and from the word model.

\section{Reflection reversal on the recurrent core}
Use indices in $\mathbb Z/k\mathbb Z$ and define
$(Qw)_i=w_{-i}$.  Then $Q^2=1$ and a direct index calculation gives
\[
 Q\rho Q=\rho^{-1}.
\]
Transporting $Q$ through $\phi$ produces an involutory reversor of the
recurrent partition permutation.  More generally, each $\rho^aQ$ is an
involutory reversor, giving $k$ phase-labelled reflection formulas; these
formulas need not act distinctly when the weight layer is nonfaithful.  The
corresponding recurrent Koopman permutation is unitary; reflection followed
by coefficientwise conjugation is an antiunitary reversor.  No such inverse
relation is claimed on all of $\mathcal P(N)$ because $T_N$ is generally
noninjective.

\section{Triangular boundary and evidence scope}
If $r=0$, the sole word is $0^k$ and the sole recurrent partition is the
staircase $(k-1,k-2,\ldots,1)$.  It is fixed, so $F_t=1$ and
$\zeta_{T_N}(z)=(1-z)^{-1}$; the recurrent Koopman spectrum is $\{1\}$,
while the full operator has zero algebraic multiplicity $p(N)-1$.  The
conclusion concerns the recurrent core; triangular deck size does not erase
the transient basin in general.

The executable census covers every $1\le N\le40$.  A producer enumerates 757
recurrent words and their 114 cycles.  A separate checker constructs all
215,307 integer partitions, follows the actual noninvertible map, and recovers
exactly the same recurrent states and cycles.  This finite agreement tests
implementations and conventions; it does not establish Brandt's all-$N$
classification or classify complete transient trees and hitting times.

The independent checker passes 658,664 assertions.  A separate SymPy path
passes 2,210 checks, producer replay is byte exact, and the hostile suite
rejects 118 semantic mutations after their payload hashes are repaired plus
one stale-hash attack.  These are implementation checks, not an independent
mathematical review.

\section{Route-A stop and limitations}
Periodic cycles and their finite determinant are exact, but partition and
necklace data contain no intrinsic rational-prime carrier, prime-power weight,
or logarithmic prime clock.  No target divisor, functional equation,
continuation, counting law, or Weil compression follows.  The finite unitary
on the recurrent core is therefore a formal operator hint only:
\[
 (A0,A1,A2,A3,A4)=(\mathrm{FAIL},\mathrm{WEAK},\mathrm{FAIL},
 \mathrm{FAIL},\mathrm{FORMAL\ HINT}).
\]
The overall verdict is \texttt{ROUTE\_A\_REJECTED}; Route B is false.  We do
not claim a complete transient tree, an exact hitting-time distribution,
nilpotent Jordan block sizes, a global reversor, or priority for the Brandt
and Akin--Davis inputs.

\paragraph{Scope.}
No target zero or prime table, arithmetic local datum, Euler factor, root
number, automorphy input, target divisor, or Hilbert--Polya operator is used.
The scope literal is \texttt{NO\_BAD\_EULER\_OR\_ROOT\_NUMBER}.

\section*{Declarations}
\small\noindent
\textbf{Data and code availability.} Package-local exact evidence and
deterministic code accompany this manuscript.
\textbf{Ethics.} No human, animal, clinical, personal, or sensitive data are
used; approval is not applicable.
\textbf{Author contributions (CRediT).} This anonymous certificate records
Conceptualization, Formal analysis, Software, Validation, and Writing; AI
systems are not authors.
\textbf{Funding and conflicts.} No external funding is reported; no conflict
is known.
\textbf{AI-use disclosure.} An AI coding assistant supported drafting and
exact-code development; it was not an external reviewer or independent error
process.\normalsize

\begin{thebibliography}{2}
\footnotesize
\raggedright
\bibitem{Brandt82} J. Brandt,
``Cycles of partitions,'' \emph{Proc. Amer. Math. Soc.} 85(3) (1982),
483--486. \url{https://doi.org/10.1090/S0002-9939-1982-0656129-5}.

\bibitem{AkinDavis85} E. Akin and M. Davis,
``Bulgarian Solitaire,'' \emph{Amer. Math. Monthly} 92(4) (1985), 237--250.
\url{https://doi.org/10.1080/00029890.1985.11971590}.
\end{thebibliography}
\end{document}
